% =====================================================================
% Regeneron STS 2027 Research Report
% Body text is VERBATIM from the URTC conference version
% (paper_current.tex, IEEEtran). Only the format changed.
%
% Format: single column, 12 pt, Times, 1.5 line spacing, 1 in margins.
%
% NOTE ON SPACING: \onehalfspacing sets Word-style 1.5 (baselinestretch
% ~1.24 at 12 pt), which is what most graders mean by "1.5 spacing."
% If STS means a literal 1.5x baseline, swap it for \setstretch{1.5}.
% Verify by measuring a rendered page: 25-26 lines/page = Word 1.5,
% 22-23 lines/page = literal 1.5.
%
% NOTE ON SECTION NUMBERING: sections are forced to Roman numerals so
% that every cross-reference resolves to the same string it did in the
% IEEEtran version (e.g. "Section V" for Discussion). Do not change
% \thesection without re-reading the prose for hard-coded numerals.
%
% NOTE ON FLOATS: the IEEEtran version combined Table I, Table II and
% Fig. 2 into a single float to keep them together in one column. That
% construct is unnecessary in a single-column layout, so they are three
% separate floats here. Declaration order is preserved, so Fig. 1-4 and
% Table I-II carry the same numbers as before.
% =====================================================================
\documentclass[12pt]{article}

\usepackage[letterpaper,margin=1in]{geometry}
\usepackage{setspace}
\usepackage{amsmath,amssymb,amsfonts}
\usepackage{newtxtext,newtxmath}     % Times text and math -- must follow amssymb
\usepackage{graphicx}
\usepackage{booktabs}
\usepackage{caption}
\usepackage{textcomp}
\usepackage{xcolor}
\usepackage{url}
\usepackage{hyperref}
\hypersetup{
  colorlinks=true,
  allcolors=black,
  urlcolor=blue,
  citecolor=black,
  linkcolor=black
}

\captionsetup{font=small,labelfont=bf,justification=justified}
\setlength{\parskip}{0pt}
\setlength{\parindent}{0.3in}

\renewcommand{\thesection}{\Roman{section}}

\title{\vspace{-0.5in}\textbf{Conformal-Gated Surrogate Screening for N-1
Under-Voltage Contingencies: Safety and Throughput Limits on a Marginally
Secure Network}}

\author{Rajan Saha\\
\small Conestoga High School, Berwyn, PA, USA\\
\small Boston University RISE Data Science Practicum}

\date{}

\begin{document}

\maketitle

\onehalfspacing

% =====================================================================
\section*{Abstract}
\addcontentsline{toc}{section}{Abstract}

The power system grid should remain resilient following the loss of any individual component. The N-1 constraint uses an AC power flow solver for each loss event of the grid elements to determine the voltage of the remaining components. With operators running every combination and repeatedly testing under different conditions, it becomes computationally expensive. Past work has either replaced these solves with surrogate models that check a group of cases at once, return a binary risk classification, or claim no missed unsafe scenarios with only a simple direct current (DC) model. In this study, we attach a one-sided split-conformal prediction interval to the predicted minimum post-contingency voltage for a ridge regression and a gradient-boosted model. For each model, there are three potential outcomes: classify the scenario as safe if the band is above 0.94 per-unit, flag it if the prediction is below the limit, or escalate the situation to the solver if it is a borderline case. On the IEEE 118-bus AC network at 90\% coverage, the gradient-boosted model is 3.29 times faster than the solver but misses 4.72\% of violations. For the gradient-boosted model, first going just under a 1\% mean missed rate at a 0.97 coverage target results in a 63.7\% escalation and speedup dropping to roughly 1.58, with error bars going slightly above 1\%. In the IEEE 30-bus system, using the thermally realistic regenerated case30 result, where the regenerated 30-bus network has 7.09\% of its contingencies in the [0.94, 0.945) boundary band compared to 56.86\% for case118, the gradient-boosted model at 0.97 coverage results in a 5.84$\pm$1.25\% escalation. These results allow operators to make an informed decision about if surrogate screening would be worthwhile on a specific network.

\vspace{0.15in}
\noindent\textbf{Keywords:} contingency screening, conformal prediction,
power-flow surrogate, selective prediction, power system security

% =====================================================================
\section{Introduction}
\label{sec:intro}

Power grids must be able to stay operational even if they lose the functionality of a single element. The N-1 criterion, for all the elements that might fail, determines whether the calculated voltages of the remaining components stay within safe operating ranges after an element fails \cite{nerc}. Normally, the check itself comes from a physics calculation called an AC power flow solve, run for each potential contingency. For a complex network with hundreds of elements, the solver becomes time-consuming and computationally expensive.

One possible solution to improve performance is to replace the slow physics solver with a faster approximation. The main issue with this approach is not that it is inaccurate on average, but when it fails, it fails undetected, declaring an unsafe contingency safe, since it only provides a single-point estimate without an attached confidence or uncertainty interval \cite{bates2021}. As a result, the operator skips the slow solver and misses the unsafe scenario altogether. The error is asymmetric, as predicting too high certifies an unsafe condition as safe while predicting too low just wastes solver time. We keep track of the solver time for escalated cases and add it to the screening gate's total time to reflect the true end-to-end cost.

There have been several approaches when it comes to employing surrogate models for N-1 contingency tests. Manoharan \cite{manoharan2026} utilized a fast computer model that checks if any single piece of equipment in the N-1 test is in danger. To double-check reliability, the system also runs randomized full solves on some of the equipment it did skip to ensure the model is accurate. Alc\'antara and Chatzivasileiadis \cite{alcantara2026} use an AI foundation model and conformal prediction to create a binary yes/no switch for post-element outages to determine whether there's a present risk or not. Christianson et al. \cite{christianson2025} designed an input-convex network that completely removes false negatives, although that claim is only relevant to DC power flow models, not the AC power flow model our study uses. Our work stands out from these papers in three distinct ways. First, unlike the first two approaches, we add a three-way gate to each contingency. Second, we specifically check for only under-voltage issues. Finally, we explain why there must be a certain floor for escalation.

% =====================================================================
\section{Background}
\label{sec:background}

A per-unit (pu) voltage normalizes power grid values into simple decimals (where 1.0 is nominal), where 0.94 per unit indicates a 6\% under-voltage. Utilities must keep voltage within a band around nominal, and this study flags any part of the grid that fails the test against an under-voltage floor of 0.94 per unit. The N-0 condition represents the current state of the grid without any equipment outages, in which all voltages are above a minimum limit. We refer to this pre-outage state of the grid as a base case. However, if a single piece of equipment fails, the rest of the network must absorb the extra load, which could cause the voltage at one or more buses to go under the safe 0.94 pu floor. N-1 screening simulates these failures and detects any violations of the under-voltage floor before they happen. In this study, we do not check for overvoltage, although 73.1\% of converged N-1 rows are over 1.05 pu, or the upper bound of the acceptable voltage band. Also, the gate only bounds one predicted scalar, so adding more factors would require their own respective target and bands. We cannot check for thermal loading because case118 assigns every line the same 9,900 MVA rating, which is a placeholder value that is not realistic, so the test would pass every case regardless of conditions.

An AC power flow solver is used to get the final voltages across a network after a single-element failure. Due to the math being complex and non-linear, it takes up to several milliseconds for one scenario, and the time multiplies across the whole grid. By contrast, a prediction from a machine learning model can take a fraction of a millisecond for predicting the single minimum voltage out of the whole grid. Classical fast screens rank power system contingencies without running a full solver \cite{ejebe1979}. While this indicates which equipment is subject to the most danger when there's an equipment failure, it fails to provide a numerical minimum voltage.

We study the IEEE 118-bus system \cite{case118}, a standard electrical power test case with 118 buses (the junctions where voltage is measured), 173 lines, and 13 transformers. Since our N-1 test set includes only lines and transformers, there are 186 total pieces of equipment that are subject to failure. We use pandapower \cite{pandapower}, a Python library used to model power systems, for solving post-outage voltages while ensuring specific generator reactive power limits are enforced.

% =====================================================================
\section{Method}
\label{sec:method}

\subsection{Dataset}

We determined the load levels by multiplying the initial network load by multipliers. For Independent mode (750 bases), we multiplied the load by multipliers ranging from 1.0 to 1.12. For Regional mode (750 bases), we did the same but from 0.9009 to 1.2310 (43.91\% outside the target range), while reactive power scales across all 1,500 bases range from 0.8156 to 1.4083 (56.83\% outside the target range). The base cases were kept if the minimum pre-outage voltage exceeded 0.94 per unit. Each accepted base case was put through a simulation of 186 contingencies representing single-element failures, followed by an alternating-current power flow calculation to determine the minimum bus voltage after each contingency. The total sample contained 1{,}500 base cases and 280{,}500 rows, including the base cases. Out of those rows, 1{,}500 were N-0 base rows that were excluded because they aren't contingencies, while the other 279{,}000 rows are N-1 contingencies. Out of those 279{,}000, 45 were removed due to failing to converge, resulting in 278{,}955 converged rows. Minimum post-contingency bus voltages ranged from 0.7179 to 0.9603 per unit, and 17.48\% of the observations violated the voltage limit. We also tested the published IEEE 30-bus system, where the published base case is at 111.83\% line loading, which is discussed in Section~\ref{sec:discussion}. Since there are 41 rather than 186 branches, the 1{,}500 bases produce 61{,}500 contingency scenarios, all of which were solved successfully.
\subsection{Surrogates}

The surrogate predicts the minimum voltage following an outage based on the per-bus loads, how generators are set up, baseline voltages, and the failed equipment. The two algorithms used in this study are a ridge linear regression model using standardized inputs and a histogram-based gradient-boosted tree ensemble \cite{sklearn}. The data is split into 60\%, 20\%, and 20\% sets using different base scenarios for training, calibration, and test sets, respectively. To tune the specific variables for each model, we split the training set again into three smaller parts: one to fit the model, one to calculate the band width $\hat{q}$, and one to test the gate on, never looking at the real calibration or test sets when choosing variable values. Instead, we keep whichever variable values skip the most solves while still missing 1\% or fewer of the violations, and we repeat the whole process with five different random splits. We test our models against two baselines: persistence, which predicts the pre-outage minimum voltage, and the training mean.

\subsection{One-sided conformal band}

The split-conformal prediction approach converts the residuals from the calibration dataset into a coverage guarantee (the number of cases the band is guaranteed to contain)\cite{lei2018,vovk2005}. Since the primary source of the danger is the voltage being less than the predicted one, we compute only the lower band. It is specified via $\hat{q}$ as the quantile for the coverage of the model overshoot (the difference between the prediction and the true value) based on the finite-sample rank approach proposed by Lei et al.\ \cite{lei2018}, and forms the one-sided band
\begin{equation}
\text{lower} = \hat{p} - \hat{q},
\end{equation}
where $\hat{p}$ is the prediction. As long as the test cases and calibration cases come from the same type of condition (a single-element outage), the true voltage remains at or above the lower bound with the desired probability. If this condition does not hold (i.e., using calibration data of a single-element outage with test data of a two-element outage), there is no guarantee for what the band truly encompasses, which raises the issue of distribution shift described in Section~\ref{sec:discussion}. A single quantile coverage is the easiest way to choose the band, and there are locally adaptive alternatives that exist \cite{romano2019}. At a 90\% coverage level, the calibrated band widths are 0.0052 per unit for the linear model and 0.0023 per unit for the gradient-boosted model.

\subsection{Three-way gate}

For each contingency, the gate performs the following tests for the band and the point prediction compared with the limit $L=0.94$:
\begin{itemize}
\itemsep0em
\item \textbf{Certify} (skip the solver): $\hat{p}-\hat{q}\ge L$, meaning that the lower edge of the band exceeds the limit.
\item \textbf{Flag} (skip the solver): $\hat{p}< L$, meaning that the point prediction itself is below the limit.
\item \textbf{Escalate} (invoke the solver) otherwise, that is, when the band straddles the limit.
\end{itemize}
The flagging test employs the point prediction rather than the lower edge of the band since the one-sided band has no upper edge (Fig.~\ref{fig:gate}). The exact solver is invoked on escalations only; the net speedup accounts for that cost:
\begin{equation}
\text{net speedup} = \frac{n\, t_{\text{solve}}}{n\, t_{\text{surr}} + n_{\text{esc}}\, t_{\text{solve}}}
\end{equation}
Here $n$ is the total number of contingencies, $t_{\text{solve}}$ is the solver time per case (set at 9.14 ms, which is the minimum over 400 timed solves), $t_{\text{surr}}$ is the surrogate time, and $n_{\text{esc}}$ is the escalation count. The net speedup is close to the inverse of the escalation rate since the surrogate is much cheaper than the solve.

% ---------------------------------------------------------------------
% FIGURE 1 (fig:gate) -- gate schematic. Cited in III-D.
% ---------------------------------------------------------------------
\begin{figure}[htbp]
\centering
\includegraphics[width=0.8\textwidth]{data/gate_schematic_v3.png}
\caption{The surrogate predicts post-outage voltage and attaches a band extending only downward. A case is certified if the whole band clears 0.94 pu, flagged if the prediction falls below it, and escalated if the band straddles it.}
\label{fig:gate}
\end{figure}

% =====================================================================
\section{Results}
\label{sec:results}

Both surrogate models accurately predict the minimum post-contingency bus voltage used by the gate to make its decision. The linear model reaches an $R^2$ of 0.77 and the gradient-boosted model 0.92, with mean absolute errors of 0.0038 and 0.0016 per unit, respectively (Table~\ref{tab:models}). At 90\% coverage, the ridge model escalates 49.1\% of cases, has 89.3\% coverage, misses 2.96\% of violations, and is 2.04 times faster than the solver. The gradient-boosted model escalates 30.6\%, has 89.8\% coverage, and is 3.29 times faster than the solver, yet it misses 4.72\% of violations. Both models ensure that the band is calibrated, as the coverage is close to 90\%. The question of safety is not in the number of errors made but in the location of those errors.

% ---------------------------------------------------------------------
% TABLE I (tab:models) -- v2 PROMOTED bodies. persistence and train-mean
% rows are copied verbatim from the committed models and are
% byte-identical.
% ---------------------------------------------------------------------
\begin{table}[htbp]
\centering
\caption{Four-model comparison at 90\% target coverage. Means $\pm$ standard deviation over five random splits. Escalation and missed-violation rate (share of true violations) in percent; histgb is the gradient-boosted model. $^{\dagger}$The train-mean model never calls the solver, so this does not represent a proper screening speedup.}
\label{tab:models}
\small
\begin{tabular}{lccccc}
\toprule
Model & MAE ($10^{-3}$ pu) & $R^2$ & Esc. & Missed & Speedup \\
\midrule
persistence & 4.2$\pm$0.1 & $-0.06\pm0.00$ & 99.5$\pm$0.2 & 0.31$\pm$0.08 & 1.00$\pm$0.00 \\
train mean  & 6.7$\pm$0.2 & $-0.00\pm0.00$ & 0.0$\pm$0.0 & 0.00$\pm$0.00 & N/A$^{\dagger}$ \\
ridge       & 3.8$\pm$0.1 & 0.77$\pm$0.01 & 49.1$\pm$2.7 & 2.96$\pm$0.44 & 2.04$\pm$0.11 \\
histgb      & 1.6$\pm$0.1 & 0.92$\pm$0.00 & 30.6$\pm$2.5 & 4.72$\pm$0.98 & 3.29$\pm$0.30 \\
\bottomrule
\end{tabular}
\end{table}

% ---------------------------------------------------------------------
% TABLE II (tab:ops) -- v2 PROMOTED (M2 gate-aware tuned) bodies, from
% notes/paper_tables_v2.tex, generated by scripts/emit_v2_tables.py from
% data/tradeoff_curve_v2.json.
% ---------------------------------------------------------------------
\begin{table}[htbp]
\centering
\caption{Safety and throughput at each coverage target on IEEE 118-bus system. Means $\pm$ standard deviation over five random splits. Escalation, empirical coverage, and missed-violation rate (share of true violations) in percent; histgb is the gradient-boosted model.}
\label{tab:ops}
\small
\begin{tabular}{llcccc}
\toprule
Model & Target & Esc. & Cov. & Missed & Speedup \\
\midrule
ridge  & 0.90 & 49.1$\pm$2.7 & 89.3$\pm$1.3 & 2.96$\pm$0.44 & 2.04$\pm$0.11 \\
ridge  & 0.94 & 64.3$\pm$2.8 & 94.0$\pm$0.7 & 0.79$\pm$0.21 & 1.56$\pm$0.07 \\
ridge  & 0.95 & 67.9$\pm$2.8 & 95.0$\pm$0.6 & 0.45$\pm$0.12 & 1.48$\pm$0.06 \\
ridge  & 0.96 & 71.4$\pm$1.4 & 95.9$\pm$0.4 & 0.14$\pm$0.07 & 1.40$\pm$0.03 \\
ridge  & 0.97 & 73.9$\pm$1.0 & 96.9$\pm$0.2 & 0.03$\pm$0.03 & 1.35$\pm$0.02 \\
ridge  & 0.98 & 74.8$\pm$0.8 & 98.0$\pm$0.2 & 0.00$\pm$0.00 & 1.34$\pm$0.01 \\
\addlinespace
histgb & 0.90 & 30.6$\pm$2.5 & 89.8$\pm$1.0 & 4.72$\pm$0.98 & 3.29$\pm$0.30 \\
histgb & 0.94 & 46.7$\pm$2.8 & 93.7$\pm$0.4 & 2.48$\pm$0.38 & 2.15$\pm$0.13 \\
histgb & 0.95 & 51.2$\pm$3.4 & 94.7$\pm$0.4 & 1.94$\pm$0.21 & 1.96$\pm$0.13 \\
histgb & 0.96 & 57.0$\pm$4.3 & 95.7$\pm$0.4 & 1.36$\pm$0.20 & 1.76$\pm$0.12 \\
histgb & 0.97 & 63.7$\pm$5.1 & 97.0$\pm$0.2 & 0.83$\pm$0.24 & 1.58$\pm$0.12 \\
histgb & 0.98 & 72.0$\pm$3.7 & 98.0$\pm$0.2 & 0.30$\pm$0.13 & 1.39$\pm$0.07 \\
\bottomrule
\end{tabular}
\end{table}

% ---------------------------------------------------------------------
% FIGURE 2 (fig:tradeoff) -- coverage sweep,
% data/tradeoff_hero_col_v2.png.
% ---------------------------------------------------------------------
\begin{figure}[htbp]
\centering
\includegraphics[width=0.65\textwidth]{data/tradeoff_hero_col_v2.png}
\caption{Escalation (solid, left axis) and missed violations (dotted, right axis) are plotted against the safety target and averaged over five splits. The mean missed rate first falls just under 1\% at 0.94 coverage for ridge and 0.97 for the gradient-boosted model, with high escalation at both and the error bars still reaching slightly above 1\%.}
\label{fig:tradeoff}
\end{figure}

\subsection{Safety comes at high escalation}

In Table~\ref{tab:ops}, we adjusted the target from 0.90 to 0.98 and recorded the missed and escalation rates. As the target gets larger, the missed rate decreases, but the escalation rate goes up because the band size has to be bigger to account for model overshoot. For the gradient-boosted model, increasing the target from 0.90 to 0.94 decreases the missed rate from 4.72\% to 2.48\%, increases the escalation from 30.6\% to 46.7\%, and decreases the speed-up from 3.29 to 2.15 times. The ridge model shows a similar pattern as well, with the average of missed cases falling below 1\% at 0.94 for the ridge model (0.79\% missed, 64.3\% escalation, 1.56 times faster) and 0.97 for the gradient-boosted model (0.83\% missed, 63.7\% escalation, 1.58 times faster), although the error bars still reach up to about 1.0--1.07\%, as Fig.~\ref{fig:tradeoff} shows.

% ---------------------------------------------------------------------
% FIGURE 3 (fig:missdepth) -- miss-depth histogram, from
% data/miss_depth_v2.png (log-y; generated by scripts/miss_depth_fig.py
% from data/missed_depth.json). Cited in IV-B.
% ---------------------------------------------------------------------
\begin{figure}[htbp]
\centering
\includegraphics[width=0.5\textwidth]{data/miss_depth_v2.png}
\caption{The histogram measures how far missed violations fall below 0.94 pu at 90\% coverage on a logarithmic count scale. Dashed lines mark each model's band width. Most misses stay within one band, but a thin tail reaches 0.0915 pu, where one certified bus had fallen to 0.8485.}
\label{fig:missdepth}
\end{figure}

\subsection{The faster model is not the safer one}

On the IEEE 118-bus system, the more accurate model is not safer at any point in the coverage axis, and Table~\ref{tab:models} shows that accuracy cannot settle this decision by itself. It can be said that the gradient-boosted model is much more accurate than the linear model, and has a 3.29 times speedup compared to 2.04 times. However, taking into account safety, it is clear that at a coverage of 0.96, the linear model misses 0.14\% of the violations while the gradient-boosted model misses 1.36\%. The linear model also shows the low missed rate of all targets during the sweep in the IEEE 118-bus system. At 0.90 coverage, 74\% of misses fall within one band width of the limit for the linear model and 55\% for the gradient-boosted model (Fig.~\ref{fig:missdepth}), with 26.7\% of ridge misses and 21.4\% of histgb misses falling more than 0.005 pu below the limit. The worst miss was a case at 0.0915 per unit under the limit, which was a bus that fell to 0.8485.

\subsection{The boundary layer sets a floor on escalation}

Most of the contingencies are at a distance of less than 0.005 per unit from the boundary, thus causing a significant number of escalations. The distribution of voltages after contingencies is shown in Fig.~\ref{fig:boundary}. Out of all the converged cases, 56.86\% are in a narrow range of [0.94, 0.945), slightly higher than the threshold, and 17.48\% lie below the limit as violations. Importantly, there is not a single 0.001-wide unit bin that accounts for more than about 14\% of the whole dataset. Certain buses account for extreme events, with bus 76 being the weakest one, holding the minimum voltage in 27.1\% of contingencies, whereas bus 53 is second with 16.81\%, followed by bus 107 with 9.31\%.

% ---------------------------------------------------------------------
% FIGURE 4 (fig:boundary) -- boundary-mass histogram, from
% data/boundary_mass_hist.png. Cited in IV-C.
% ---------------------------------------------------------------------
\begin{figure}[htbp]
\centering
\includegraphics[width=0.8\textwidth]{data/boundary_mass_hist.png}
\caption{Lowest voltage across all failure cases; anything left under the line is a violation. Over half are within 0.005 pu above the limit, while the tail below the limit drops much lower.}
\label{fig:boundary}
\end{figure}

A case is escalated when the prediction lies within one band-width above the 0.94 pu limit, since the band straddles the threshold. The gradient-boosted model's band is only 0.0023 pu wide, but so much of the distribution sits just above 0.94 that this narrow window still captures 30.6\% of all contingencies. As the band becomes larger, escalation moves closer to every case predicted at or above the limit, the remaining 82.52\% of cases, or the saturation point. Linear prediction overpredicts cases below the limit compared to gradient-boosted prediction, and thus has a lower escalation ceiling of 74.89\% versus 82.79\%. The gradient-boosted model flags violations about as often as they occur, so its ceiling essentially lands on the saturation point. A model that does not frequently predict below the threshold has a higher ceiling and misses more violations.

Persistence is the baseline model where the voltage before the outage of one element is assumed to be equal to the voltage after it. All the pre-outage scenarios are above 0.94 pu, so its predictions never go below the limit, and since it cannot predict under the limit, it can never mark a case as a violation. In Table~\ref{tab:models}, the model escalates 99.5\% of cases and is barely faster than the solver in terms of time taken. In order to reduce the rate of escalation, the model will need to confirm that there are safe cases, but it can also risk misclassifying some of them. An operator that requires a missed rate around 1\% (with error bars still touching the 1\% threshold) would run the models at a 0.94 or a 0.97 target coverage and accept about a 1.6 times speedup instead of the 2 to 3.3 times available at 0.90.

% =====================================================================
\section{Discussion and limitations}
\label{sec:discussion}

Calibrated uncertainty, as well as making the decision of either outsourcing the problem or not to a more complex model \cite{desalvo2015} and providing the model with multiple choices \cite{angelopoulos2024,cortes2016,chow1970}, are proven methods. In our study, two things differ from that literature. First, we fall back to an exact solver instead of a more complex model, so the escalated case can be resolved, and second, we also explain why the speedup is capped, which is because most of the data is very close to the limit.

Our approach differs from Manoharan's audited population control \cite{manoharan2026}, which is itself a form of distribution-free risk control \cite{bates2021}. His method uses a bound for the proportion of skipped cases that are unsafe, while our gate uses a margin of error that is subtracted from each prediction. This establishes the lower bound because, in split conformal prediction, we can be correct on average but not necessarily inside the band. If one demands perfect correctness with a limited amount of calibration data, then the prediction intervals become extremely wide, to the point where it becomes useless \cite{barber2021}. Every base case generates 186 related contingencies, and therefore, the coverage rates are measured averages and not an assurance for a single contingency.

The key findings obtained during the experiments relate to a network with a realistic topology and synthetic loads and generation. This escalation floor is set by how the distribution of the data occurs in relation to the 0.94 pu limit. For instance, in a thermally realistic regeneration of the IEEE 30-bus network, 7.09\% of contingencies fall just above the 0.94 pu limit, while another 15.40\% fall below it. For case118, 56.86\% fall just above the limit, and 17.48\% fall below it. The original 30-bus system is thermally unrealistic, since it has 111.83\% line loading in its base case. Using the thermally realistic regeneration, the gradient-boosted model with a coverage target of 0.97 gives 5.84$\pm$1.25\% escalations. The target for which the miss rate becomes less than 1\% remains unresolved after five split-ups: for 0.96 it is 0.91$\pm$0.22\%, with only two of five being under this limit. We use 0.94 pu as a conservative choice. Only 86 of the 1,500 base cases are above 0.95 pu, and increasing the threshold to 0.95 pu leads to a ridge escalation of 1.38$\pm$0.33\%. In other words, it finds everything hazardous.

All metrics reported use the sampled N-1 population, which includes approximately one-quarter of cases involving both a branch outage and a single-generator outage. Moreover, mathematical proof is not capable of guaranteeing the safety of N-2 outages because the use of conformal prediction assumes that the data is exchangeable with the calibration set \cite{vovk2005}. Also, we only checked if the final voltage was too low, not for thermal overheating, over-voltage, or whether or not the grid crashed in the first place, which is called dynamic stability. The band only counts instances where the model is wrong and not how badly it fails. In addition, the worst case, at 0.0915 per unit, represents a case of sudden collapse where the generator at that particular worst-case scenario's weakest bus reaches the limit of reactive power and can no longer control the voltage. The model cannot foresee this event just from the pre-outage data. Lastly, all of the testing data were obtained by making slight modifications to the same pre-contingency grid setup. Therefore, there is some uncertainty as to how often the models should be recalibrated due to changes in the seasons and changes in demand.

% =====================================================================
\section{Conclusion and future work}
\label{sec:conclusion}

Our models can predict N-1 under-voltage contingencies on the IEEE 118-bus system. On this network, safety and speed are inversely related, though the 30-bus result shows that this depends on data distribution relative to the limit. At 90\% coverage, both methods perform very fast, yet they miss too many violations to be considered safe. To make sure that the missed ratio approaches 1\% on average, almost two-thirds of the cases have to be escalated, resulting in a speed-up factor of 1.6. The main reason is that most of the cases fall very close to the limit. Using the thermally realistic regeneration, the bands often straddle this limit and then must escalate to the solver as the result is uncertain. The novel element here is the reasoning for the existence of such a floor, and a warning that people should expect less speed-up from a surrogate model that is calibrated to be safe on average over the N-1 distribution. Future work involves examining multi-element outages and a third network to determine whether speedup depends on proximity to the threshold.


\section*{Acknowledgments}

This paper has been written as part of the Boston University RISE data science practicum. The authors would like to thank the program and the teaching fellows and staff. An AI assistant, Claude from Anthropic, was used to validate the code and grammar.
% =====================================================================
% References. arXiv entries verified in notes/prior-art.md; the classic and
% software entries were fetch-verified on 2026-07-21 (dblp, IEEE Xplore,
% Oxford Academic, NeurIPS proceedings, JMLR, UW test-case archive); the two
% ALT conference entries (cortes2016, desalvo2015) on 2026-07-23 (Crossref + dblp;
% Springer chapter pages behind an auth wall); bates2021 from the publisher
% PDF; tibshirani2019 on 2026-07-26 (NeurIPS proceedings page + ACM DL);
% angelopoulos2024, christianson2025 and ansi2020 on 2026-07-29 (medRxiv full
% text, proceedings.mlr.press/v283/christianson25a.html, and the NEMA official
% excerpt respectively). No field is filled from memory.
% NERC (nerc) confirmed against nerc.com as TPL-001-5.1, effective
% 2023-07-01 -- ADD THE VERIFICATION DATE HERE, every other entry has one.
% =====================================================================
\singlespacing
\begin{thebibliography}{17}

\bibitem{nerc} \textit{Transmission System Planning Performance Requirements}, NERC Reliability Standard TPL-001-5.1, North American Electric Reliability Corporation, Atlanta, GA, USA, eff. Jul. 1, 2023.

\bibitem{bates2021} S. Bates, A. Angelopoulos, L. Lei, J. Malik, and M. Jordan, ``Distribution-free, risk-controlling prediction sets,'' \textit{J. ACM}, vol. 68, no. 6, art. 43, 2021.

\bibitem{manoharan2026} J. Manoharan, ``Audited selective verification for risk-controlled N-1 thermal contingency screening under deployment shift,'' arXiv:2607.13221, 2026.

\bibitem{alcantara2026} A. Alc\'antara and S. Chatzivasileiadis, ``Trustworthiness layer for foundation models in power systems: Application to N-k contingency screening,'' arXiv:2602.07995v2, 2026.

\bibitem{christianson2025} N. Christianson, W. Cui, S. Low, W. Yang, and B. Zhang, ``Fast and reliable N-k contingency screening with input-convex neural networks,'' in \textit{Proc. 7th Annu. Learn. Dynamics Control Conf. (L4DC)}, ser. PMLR, vol. 283, 2025, pp. 527--539.

\bibitem{ejebe1979} G. C. Ejebe and B. F. Wollenberg, ``Automatic contingency selection,'' \textit{IEEE Trans. Power App. Syst.}, vol. PAS-98, no. 1, pp. 97--109, 1979.

\bibitem{case118} R. Christie, ``Power systems test case archive: IEEE 118-bus system,'' Univ. of Washington, 1993. [Online]. Available: https://labs.ece.uw.edu/pstca/

\bibitem{pandapower} L. Thurner, A. Scheidler, F. Sch\"afer, \textit{et al.}, ``pandapower---an open-source Python tool for convenient modeling, analysis, and optimization of electric power systems,'' \textit{IEEE Trans. Power Syst.}, vol. 33, no. 6, pp. 6510--6521, Nov. 2018.

\bibitem{sklearn} F. Pedregosa \textit{et al.}, ``Scikit-learn: Machine learning in Python,'' \textit{J. Mach. Learn. Res.}, vol. 12, pp. 2825--2830, 2011.

\bibitem{lei2018} J. Lei, M. G'Sell, A. Rinaldo, R. J. Tibshirani, and L. Wasserman, ``Distribution-free predictive inference for regression,'' \textit{J. Amer. Statist. Assoc.}, vol. 113, no. 523, pp. 1094--1111, 2018.

\bibitem{vovk2005} V. Vovk, A. Gammerman, and G. Shafer, \textit{Algorithmic Learning in a Random World}. New York, NY, USA: Springer, 2005.

\bibitem{romano2019} Y. Romano, E. Patterson, and E. J. Cand\`es, ``Conformalized quantile regression,'' in \textit{Adv. Neural Inf. Process. Syst. (NeurIPS)}, vol. 32, 2019, pp. 3538--3548.

\bibitem{barber2021} R. F. Barber, E. J. Cand\`es, A. Ramdas, and R. J. Tibshirani, ``The limits of distribution-free conditional predictive inference,'' \textit{Inf. Inference: J. IMA}, vol. 10, no. 2, pp. 455--482, 2021.

\bibitem{desalvo2015} G. DeSalvo, M. Mohri, and U. Syed, ``Learning with deep cascades,'' in \textit{Algorithmic Learning Theory (ALT)}, ser. Lecture Notes in Comput. Sci., vol. 9355. Springer, 2015, pp. 254--269.

\bibitem{angelopoulos2024} A. N. Angelopoulos \textit{et al.}, ``Conformal triage for medical imaging AI deployment,'' medRxiv, 2024, doi: 10.1101/2024.02.09.24302543.

\bibitem{cortes2016} C. Cortes, G. DeSalvo, and M. Mohri, ``Learning with rejection,'' in \textit{Algorithmic Learning Theory (ALT)}, ser. Lecture Notes in Comput. Sci., vol. 9925. Springer, 2016, pp. 67--82.

\bibitem{chow1970} C. K. Chow, ``On optimum recognition error and reject tradeoff,'' \textit{IEEE Trans. Inf. Theory}, vol. 16, no. 1, pp. 41--46, 1970.

\end{thebibliography}

\end{document}